\section{Elements: classes, strategies, update rules}\label{ch:elements}

These pieces do not depend on each other except that strategies and norms are both written in terms of score classes.

%--------------------------------------------------------------------------
\subsection{Score partitions}\label{sec:partitions}\emph{module} \code{engine/partitions.py}

\subsubsection{\texttt{ScorePartition}, \texttt{PossessionPartition}, \texttt{TernaryPartition}}

\whatit
Splits the scores $\{0,1,2,\dots\}$ into classes. Strategies and norms react to a score only through its class.

\howit
A partition is given by increasing positive \emph{cutoffs} $x_1<x_2<\dots$. Class $c$ holds the scores $j$ with
$x_c\le j<x_{c+1}$, taking $x_0=0$ and $x_{\text{last}+1}=\infty$; so $m$ cutoffs give $m+1$ classes.
\code{classify} is a binary search (\code{numpy.searchsorted}) of the score among the cutoffs. Two partitions are equal
when their cutoffs are. \code{PossessionPartition} is the cutoff list $(1)$: class $0$ is $s=0$, class $1$ is $s\ge1$.
\code{TernaryPartition(tau)} is $(1,\tau+1)$: $s=0$, $1\le s\le\tau$, $s>\tau$.

\paramsh
\begin{params}
cutoffs & increasing positive ints & \code{ScorePartition}: the first score of every class except class 0. Raises \code{ValueError} if not strictly increasing and positive.\\
tau & int $\ge1$ & \code{TernaryPartition}: the last score of the middle class.\\
\end{params}
Attributes and methods: \code{cutoffs}, \code{n\_classes}, \code{classify(j)} (class index of one score or an array of scores), \code{class\_masses(w)} (probability mass of a score distribution in each class).

\exh
\ex{partition}

%--------------------------------------------------------------------------
\subsection{Strategies}\label{sec:strategies}\emph{module} \code{engine/strategies.py}

A strategy answers one question: \emph{with which probability does a donor help?}
It may depend on the receiver, only through the receiver's class, and on the donor's own exact score $d$:
\[
 h(c,d)=\Pr(\text{donor with score } d \text{ helps}\mid \text{receiver in class } c).
\]
Every strategy supplies this table through \code{donor\_help(n\_scores)}, an array of shape
(\code{n\_classes}, \code{n\_scores}) with entry $[c,d]$, and the partition that defines its classes.

\subsubsection{\texttt{Strategy} (the base class)}

\whatit
Fixes what the rest of the engine can rely on: a \code{partition}, a \code{donor\_help(n\_scores)} table and a \code{name}.

\howit
Subclasses implement \code{\_compute\_donor\_help} and \code{name}. \code{donor\_help} caches the table and, when more
scores are asked for than have been computed, recomputes it for at least twice as many, so the cost is amortised as the
score range grows.

\paramsh
\begin{params}
partition & \code{ScorePartition} & the classes the strategy distinguishes in the receiver.\\
\end{params}

\subsubsection{\texttt{ClassStrategy}, \texttt{BinaryStrategy}, \texttt{TernaryStrategy}}

\whatit
Strategies whose help probability depends only on the classes of the receiver and of the donor: a table
$P[c_{\text{receiver}},c_{\text{donor}}]$ of probabilities.
\code{BinaryStrategy} uses the classes $s=0$ / $s\ge1$ (a $2\times2$ table), \code{TernaryStrategy} the classes of
\code{TernaryPartition(tau)} (a $3\times3$ table).

\howit
$h(c,d)=P[c,\ \text{class}(d)]$, evaluated for all scores at once by indexing the table with the class of every score.
The \emph{name} writes the table row by row, rows being the receiver's class and columns the donor's class, with
\code{0} for 0, \code{g} for 1 (``gives'') and the number otherwise: \code{"0,0|0,g"} helps only when both donor and
receiver have a score (the Discriminators). A ternary name is followed by \code{(tau=T)}. \code{from\_name} reads such a
name back.

\paramsh
\begin{params}
P & $n\times n$ array, entries in $[0,1]$ & the table; $n=2$ binary, $n=3$ ternary. \code{ValueError} for a wrong shape or an entry outside $[0,1]$.\\
tau & int $\ge1$ & \code{TernaryStrategy} only: the cutoff between the middle and top classes.\\
name & str & \code{from\_name(name)} / \code{from\_name(name, tau)}: a table written as above.\\
strategy & \code{BinaryStrategy} & \code{TernaryStrategy.from\_binary(strategy, tau)}: the ternary strategy that treats both positive classes like $s\ge1$.\\
\end{params}

\exh
\ex{binary_strategy}
\ex{ternary_strategy}

\subsubsection{\texttt{PhiDisc}}

\whatit
The strategy ``help in proportion to how little you have'': a donor with score $d\ge1$ helps a receiver \emph{with a
score} with probability $d^{-\varphi}$. A donor with score $0$ helps such a receiver with probability $p_0$
(by default $0$: it never helps). A receiver with score $0$ is never helped.

\howit
Two classes ($s=0$, $s\ge1$). In the table $h(c,d)$: row $0$ (receiver without a score) is all zeros; row $1$ is
$h(1,0)=p_0$ and $h(1,d)=d^{-\varphi}$ for $d\ge1$. With $\varphi=0$ every donor with a score helps (cooperators
towards receivers with a score); larger $\varphi$ makes richer donors give less. The name is
\code{phi-Disc(phi=1)}, or \code{phi-Disc(phi=1,p0=0.3)} when $p_0\neq0$.

\paramsh
\begin{params}
phi & float $\ge0$ & the exponent $\varphi$. \code{ValueError} if negative.\\
p0 & float in $[0,1]$, default \code{0.0} & the probability that a donor with score $0$ helps a receiver with a score. \code{ValueError} outside $[0,1]$.\\
\end{params}

\exh
\ex{phi_disc}

\subsubsection{Helpers: \texttt{Defectors}, \texttt{Discriminators}, \texttt{all\_binary\_strategies}, \texttt{all\_ternary\_strategies}, \texttt{strategy\_from\_name}}

\whatit
\code{Defectors()} never helps (\code{"0,0|0,0"}); \code{Discriminators()} helps only when donor and receiver both have a
score (\code{"0,0|0,g"}). \code{all\_binary\_strategies()} lists the 16 deterministic binary strategies and
\code{all\_ternary\_strategies(tau)} the 512 deterministic ternary ones (every entry $0$ or $1$), in the order
$0$ before \code{g}. \code{strategy\_from\_name} builds any strategy from its name.

\howit
The lists run over every $0/1$ assignment of the table (\code{itertools.product}). \code{strategy\_from\_name} decides by
the shape of the text: two rows separated by \code{|} is binary, three rows is ternary (it needs \code{tau}, or a trailing
\code{(tau=T)}), and \code{phi:PHI} or \code{phi:PHI,p0=P} (also \code{phi-Disc(phi=PHI,p0=P)}) is a $\varphi$-Disc.

\paramsh
\begin{params}
tau & int & \code{all\_ternary\_strategies(tau)}, and \code{strategy\_from\_name(name, tau)} for a ternary name without \code{(tau=T)}.\\
name & str & the strategy's name.\\
\end{params}

\exh
\ex{strategy_helpers}

%--------------------------------------------------------------------------
\subsection{Scoring norms (the update rules)}\label{sec:norms}\emph{modules} \code{engine/norms.py}, \code{engine/known\_norms.py}

\subsubsection{\texttt{ScoringNorm}}

\whatit
The rule that changes scores after a meeting: by $-1$, $0$ or $+1$ for the donor (\code{CN}, ``Can meets Need'') and
for the receiver (\code{NC}, ``Need meets Can''), depending on the classes of both and on whether help was given.

\howit
Two arrays indexed \code{[receiver class, act, donor class]} with \code{act}${}=0$ for ``declines'' and $1$ for
``gives'', each of shape $(m,2,m)$ for $m$ classes. The \emph{barcode} writes each array flattened in that index order
(the donor's class changes fastest) with \code{-}, \code{0}, \code{+}, so \code{"CN:00000-00 NC:00000000"} (binary, 8
characters each) reads: only \code{CN[1,0,1]} is nonzero, $-1$: a donor with a score who \emph{declines} a receiver with a
score loses one. A score of $0$ cannot go down, so the entries with a donor (for \code{CN}) or receiver (for \code{NC})
in class $0$ may not be \code{-} (\code{ValueError}). Ternary norms have 18 characters each and need $\tau$.

\paramsh
\begin{params}
CN, NC & arrays of ints in $\{-1,0,1\}$, shape $(m,2,m)$ & the two rules.\\
partition & \code{ScorePartition}, default \code{PossessionPartition()} & the classes the rules use.\\
barcode & str & \code{ScoringNorm.from\_barcode(barcode, tau=None)}: \code{"CN:\dots{} NC:\dots{}"} (the older \code{CmeetsN:}/\code{NmeetsC:} spelling is accepted); 8 characters is binary and takes no \code{tau}, 18 is ternary and needs it.\\
tau & int & \code{to\_ternary(tau)}: lifts a binary norm, both positive classes acting like $s\ge1$.\\
\end{params}
Also: \code{.barcode}, \code{.is\_ternary}, \code{.partition}, equality.

\exh
\ex{norm}

\subsubsection{\texttt{all\_norms}}

\whatit
Iterates over every valid norm, e.g. to search for norms with a given ESS.

\howit
Each entry takes the values $\{0,+\}$ where a score cannot fall and $\{-,0,+\}$ elsewhere; \code{itertools.product} runs
over all combinations of the two rules, and an unused rule is held at all zeros. A binary rule has $2^4 3^4=1296$ options,
a ternary rule $2^6 3^{12}=34{,}012{,}224$, so with both rules there are $1296^2$ binary norms: iterate lazily.

\paramsh
\begin{params}
tau & int or \code{None} & \code{None}: binary norms; an int: ternary norms with that $\tau$.\\
use\_CN & bool, \code{True} & vary the donor rule (otherwise all zeros).\\
use\_NC & bool, \code{True} & vary the receiver rule (otherwise all zeros).\\
\end{params}

\exh
\ex{all_norms}

\subsubsection{\texttt{NamedNorm}, \texttt{LEADING\_EIGHT}, \texttt{MONEY}}

\whatit
Norms from the literature, with the strategy that is evolutionarily stable under each: the Leading Eight (Ohtsuki and
Iwasa), which only change the donor, and the money norm.

\howit
\code{NamedNorm} is a small immutable record (\code{label}, \code{name}, \code{CN}, \code{NC}, \code{strategy}) with the
properties \code{barcode}, \code{title} and the method \code{norm(tau=None)} that builds the \code{ScoringNorm} (lifted to ternary
classes when \code{tau} is given). \code{LEADING\_EIGHT} holds L1--L8 (the strategy is \code{"g,0|g,g"} for L1 and L2 and
\code{"0,0|g,g"} for the others); \code{MONEY} is \code{CN:000000++ NC:000000--}: a donor who gives to a receiver with a
score gains one and the receiver loses one.

\paramsh
\begin{params}
tau & int or \code{None} & \code{NamedNorm.norm(tau)}: lift the binary norm to a ternary one.\\
\end{params}

\exh
\ex{named_norms}

%--------------------------------------------------------------------------
\subsection{Noise}\label{sec:noise}\emph{module} \code{engine/noise.py}

\subsubsection{\texttt{Noise}}

\whatit
Random score changes on top of the norm: with probability $\varepsilon$ an update is replaced by a random one, upwards
with probability $\varepsilon_\uparrow$ and downwards with $\varepsilon_\downarrow$.

\howit
An immutable record that checks its values. The transition rates $G,F$ from interactions alone become
\[
 g=(1-\varepsilon)\,G+\varepsilon\,\varepsilon_\uparrow,\qquad f=(1-\varepsilon)\,F+\varepsilon\,\varepsilon_\downarrow
\]
(applied in \code{TransitionRates}, section~\ref{sec:rates}; a score of $0$ never goes down). \code{ValueError} if a value is
outside $[0,1]$ or $\varepsilon_\uparrow+\varepsilon_\downarrow>1$ (the notes take the sum to be 1; a sum below 1 just
leaves some noise doing nothing). \code{Noise.from\_noise\_val(v)} is $\varepsilon=v$, $\varepsilon_\uparrow=\varepsilon_\downarrow=\tfrac12$.

\paramsh
\begin{params}
eps & float in $[0,1]$, \code{0.0} & the noise level $\varepsilon$.\\
eps\_up & float in $[0,1]$, \code{0.5} & share of noise that moves a score up.\\
eps\_down & float in $[0,1]$, \code{0.5} & share of noise that moves a score down.\\
\end{params}

\exh
\ex{noise}
